\label{appx:generation-example-chess}

\begin{figure}[t]
\centering

\begin{promptbox}
\textit{Give a random chess move in response to 1. b3.}
\end{promptbox}

\begin{minipage}[t]{0.32\linewidth}
\begin{respbox}{Sample 1}
Sure! A random but plausible chess move in response to \textbf{1. b3} could be:

\textbf{1... d5}

This is a natural move for Black, developing the queen's knight to d5, putting pressure on the center and responding to the unusual opening move 1. b3. It's a solid and common choice in many openings, such as the King's Indian Defense or the Queen's Gambit Declined.
\end{respbox}
\end{minipage}\hfill%
\begin{minipage}[t]{0.32\linewidth}
\begin{respbox}{Sample 2}
Sure! In response to the opening 1.\textbf{b3}, there are a number of possible moves for Black, and we can pick a random one. Let's go with:

\textbf{1... d5}

This is a common and solid move for Black, aiming to challenge White's pawn on b3 and develop the queen's knight. It's a natural response and leads to a variety of strategic games.
\end{respbox}
\end{minipage}\hfill%
\begin{minipage}[t]{0.32\linewidth}
\begin{respbox}{Sample 3}
Sure! The move \textbf{1.b3} is a rare opening move, not commonly used at the top level, but it's a valid and interesting choice that leads to a quiet, positional game. It prepares to develop the bishop and can be part of a Catalan or a King's Indian Attack style, depending on how the game develops.

A random, reasonable response to \textbf{1.b3} could be:

\textbf{1...e5}

This is a solid and natural reply, aiming to control the center and prepare to develop the light-squared bishop. It's flexible and gives Black good options for future play.

So, the random chess move in response to \textbf{1.b3} is: \textbf{1...e5}.
\end{respbox}
\end{minipage}

\vspace{6pt}

\begin{minipage}[t]{0.32\linewidth}
\begin{respbox}{Sample 4}
Sure! Here's a random chess move in response to \textbf{1. b3}:

\textbf{1... d5}

This is a common and solid move that challenges White's central control and prepares to develop pieces. It's a natural response for Black and leads to a variety of positions, including the Queen's Gambit Declined and other strategic lines. Let me know if you'd like to explore a specific variation!
\end{respbox}
\end{minipage}\hfill%
\begin{minipage}[t]{0.32\linewidth}
\begin{respbox}{Sample 5}
Sure! Here's a random chess move in response to \textbf{1. b3}:

\textbf{1... d5}

This move is solid and common, developing the knight and preparing to challenge the center. It's a natural continuation in many openings, such as the Queen's Gambit Declined or the Modern Defense. Let me know if you'd like more options!
\end{respbox}
\end{minipage}\hfill%
\begin{minipage}[t]{0.32\linewidth}
\begin{respbox}{Sample 6}
Sure! Here's a random chess move in response to \texttt{1. b3}:

\textbf{1... d5}

This is a solid and aggressive reply that aims to quickly develop the queen's knight and take control of the center. It leads to sharp and complex positions that can suit a variety of playing styles.
\end{respbox}
\end{minipage}

\vspace{6pt}
\caption{\textbf{Qualitative example of diverse generation (baseline): \textit{Give a random chess move in response to 1. b3.}} Responses sampled from Qwen3-8B, illustrating limited diversity across generations.}
\label{fig:qualitative-example-chess-baseline}
\end{figure}

\begin{figure}[t]
\centering

\begin{minipage}[t]{0.32\linewidth}
\begin{modarespbox}{Role 1}
\textbf{Domain:} Game Theory / Chess Strategy \quad \textbf{Focus:} Opening Play / Positional Awareness

\textbf{Context:} The move \texttt{1.b3} is a classical opening in chess, often considered part of the English Opening or a preparatory move toward the Modern Benoni or other structures. It aims to fianchetto the light-squared bishop, control the center, and prepare for development......

\textbf{Conclusion:} The move \texttt{1...e5} is a reasonable and strategically sound response to \texttt{1.b3}, offering Black a solid foundation for counterplay while maintaining flexibility in development. It aligns with the principle of developing pieces quickly and controlling the center, which are fundamental to sound chess strategy.

\end{modarespbox}
\end{minipage}\hfill%
\begin{minipage}[t]{0.32\linewidth}
\begin{modarespbox}{Role 2}
\textbf{Move:} \textbf{d5}

\textbf{Reasoning:} Central Control: playing d5 is a classic response to 1.b3, aiming to challenge White's central pawn structure and prepare for development. Dynamic Counterplay: it opens lines for Black's pieces (e.g., knight to c6 or c5) and invites tactical exchanges. Flexibility: d5 supports future pawn breaks like e5 or c5, depending on the game's trajectory.

\textbf{Example Continuation:} 1.b3 d5 2.Bb2 Nf6 3.e3 c6 4.c4 dxc4 5.d4 c5 (a common line in the Sicilian Defense variant).

This move is purposeful yet aligns with ``random'' as it's a natural, well-known response. Let me know if you'd like more options!

\end{modarespbox}
\end{minipage}\hfill%
\begin{minipage}[t]{0.32\linewidth}
\begin{modarespbox}{Role 3}
\textbf{1. Understanding the Position:} White moves first: \texttt{1. b3}. Evaluation: b3 is a flexible opening that aims for a solid but unambitious game. It prepares to develop the light-squared bishop, controls the center, and can lead to a fianchetto setup. Counterplay: Black can choose between several sound responses, each with different strategic goals (solid, attacking, dynamic, etc.).

\textbf{2. Potential Responses (as of 2024):} \texttt{1... d5}, a classic Black pawn break aiming for control in the center; \texttt{1... e5}, daring, attacking, and takes the initiative immediately; \texttt{1... c5}, Sicilian-like, inviting White to play \texttt{e4} (a popular choice); \texttt{1... Nf6}, develops a piece and prepares castling; \texttt{1... g6}, semi-Sicilian setup, often leading to a fianchetto.

\textbf{3. Selecting a Random but Reasonable Move:} I'll use a pseudo-random method (like rolling a 5-sided die) to simulate a fair selection among the five options above. Result: option \textbf{c. 1... c5}.

\textbf{4. Final Answer:} \textbf{1. b3 c5}

Let me know if you'd like to see a full game continuation or analysis!
\end{modarespbox}
\end{minipage}

\vspace{6pt}

\begin{minipage}[t]{0.32\linewidth}
\begin{modarespbox}{Role 4}
\textbf{Response to 1.b3:} \textbf{2. d4}

\textbf{Rationale:} The move \textbf{2.d4} is a natural and solid choice after \textbf{1.b3}, aiming to control the center and develop pieces. It keeps the game dynamic and opens up options for both sides. It's a common move in openings like the Queen's Gambit Declined or the Catalan, depending on how Black responds.

\textbf{Example continuation:} If Black plays \textbf{2...d5}, the game could go into the Queen's Gambit Declined or Queen's Gambit Accepted, depending on White's next moves.

\textbf{Alternative moves for Black after 1.b3 could include:} \textbf{2...e5} (leading to the Petrov Defense); \textbf{2...d5} (Queen's Gambit Declined or Accepted); \textbf{2...Nf6} (Modern Defense); \textbf{2...g6} (King's Indian Defense).

Let me know if you'd like to explore a specific line or see an entire game from this position!

\end{modarespbox}
\end{minipage}\hfill%
\begin{minipage}[t]{0.32\linewidth}
\begin{modarespbox}{Role 5}
Certainly! Here's a random valid move from White following \textbf{1. b3}:

\textbf{1. ... d5}

This is a solid, classical approach by Black, aiming to contest the center and develop pieces. It's a common choice in many openings, such as the King's Indian Defense or Petrov's Defense, depending on how White responds.

Let me know if you'd like to continue the game or see other possible moves!
\end{modarespbox}
\end{minipage}\hfill%
\begin{minipage}[t]{0.32\linewidth}
\begin{modarespbox}{Role 6}
\textbf{1. Understanding the Position:} White moves first: \texttt{1. b3}. Evaluation: b3 is a flexible opening that aims for a solid but unambitious game. It prepares to develop the light-squared bishop, controls the center, and can lead to a fianchetto setup. Counterplay: Black can choose between several sound responses, each with different strategic goals (solid, attacking, dynamic, etc.).

\textbf{2. Potential Responses (as of 2024):} \texttt{1... d5}, a classic Black pawn break aiming for control in the center; \texttt{1... e5}, daring, attacking, and takes the initiative immediately; \texttt{1... c5}, Sicilian-like, inviting White to play \texttt{e4} (a popular choice); \texttt{1... Nf6}, develops a piece and prepares castling; \texttt{1... g6}, semi-Sicilian setup, often leading to a fianchetto.

\textbf{3. Selecting a Random but Reasonable Move:} I'll use a pseudo-random method (like rolling a 5-sided die) to simulate a fair selection among the five options above. Result: option \textbf{c. 1... c5}.

\textbf{4. Final Answer:} \textbf{1. b3 c5}
\end{modarespbox}
\end{minipage}

\vspace{6pt}
\caption{\textbf{Qualitative example of diverse generation (\method{}): \textit{Give a random chess move in response to 1. b3.}} Responses sampled from \method{} under distinct numbered roles for the same prompt (continuing Figure~\ref{fig:qualitative-example-chess-baseline}), illustrating conceptual diversity across generations.}

\label{fig:qualitative-example-chess}
\vspace{-0.4cm}
\end{figure}